% GENERATED FILE. Edit canonical lessons, then run npm run generate:books.
% canonical-source-hash: fnv1a64:74ee7669e2de85be
% canonical-lessons: FR-C285-rumeur, FR-C285-scoop, FR-C285-desinformation, FR-C285-censure, FR-C285-polemique

\chapter{On dit que…}
\label{ch:fr-on-dit-que}
\hlchaptermodality{285}

\begin{chapteropening}
\textbf{By the end of this chapter:} \emph{I can talk about what goes round: a rumour, a scoop, false information, censorship and a controversy.}

Complete the last lesson of chapter 285: I can talk about what goes round: a rumour, a scoop, false information, censorship and a controversy.
\end{chapteropening}

\section[la rumeur]{la rumeur — the rumour}

\label{lesson:FR-C285-rumeur}



\begin{quote}
Before the new one: say the French for to recount, then the French for to comment on.
\end{quote}

\subsection*{You'll want to know: la rumeur}
\textbf{la rumeur} — \textquotedblleft{}the rumour\textquotedblright{}.

\begin{quote}
Une rumeur court dans le village.
\end{quote}

\begin{quote}
\emph{A rumour is going round the village.}
\end{quote}

\emph{La rumeur court} — the rumour is going round. Also a low noise: \emph{la rumeur de la ville}.

\begin{cousinweb}
\emph{rumeur} is inherited from Latin \emph{rūmor} \textquotedblleft{}common talk, hearsay\textquotedblright{} (EtymDB).
\end{cousinweb}

\begin{grammarlens}[title={What you've built}]
One more word for telling, explaining and describing at B1.
\end{grammarlens}

\subsection*{Your turn}
\begin{itemize}
\raggedright
  \item \emph{Say it:} \emph{la rumeur}
  \item \emph{Say it:} \emph{la rumeur}, once more
  \item \emph{Say it:} a sentence of your own with \emph{la rumeur}
  \item \emph{Recall:} say the French for to recount, then the French for to comment on, then say \emph{la rumeur} again
\end{itemize}

\begin{tcolorbox}[breakable,colback=teal!4,colframe=teal!35!black,title={Before you move on}]
What does la rumeur mean? (\textquotedblleft{}The rumour\textquotedblright{}.) Say it once more.
\end{tcolorbox}
\section[le scoop]{le scoop — the scoop, the exclusive story}

\label{lesson:FR-C285-scoop}



\begin{quote}
Before the new one: say the French for to comment on, then the French for the rumour.
\end{quote}

\subsection*{You'll want to know: le scoop}
\textbf{le scoop} — \textquotedblleft{}the scoop, the exclusive story\textquotedblright{}.

\begin{quote}
Ce journal a eu le scoop avant tout le monde.
\end{quote}

\begin{quote}
\emph{This paper got the scoop before everyone else.}
\end{quote}

Borrowed from English; official French proposes \emph{l'exclusivité}.

\begin{cousinweb}
\emph{scoop} was borrowed from English \emph{scoop} \textquotedblleft{}a ladle; a story got ahead of rivals\textquotedblright{} (EtymDB).
\end{cousinweb}

\begin{grammarlens}[title={What you've built}]
One more word for telling, explaining and describing at B1.
\end{grammarlens}

\subsection*{Your turn}
\begin{itemize}
\raggedright
  \item \emph{Say it:} \emph{le scoop}
  \item \emph{Say it:} \emph{le scoop}, once more
  \item \emph{Say it:} a sentence of your own with \emph{le scoop}
  \item \emph{Recall:} say the French for to comment on, then the French for the rumour, then say \emph{le scoop} again
\end{itemize}

\begin{tcolorbox}[breakable,colback=teal!4,colframe=teal!35!black,title={Before you move on}]
What does le scoop mean? (\textquotedblleft{}The scoop, the exclusive story\textquotedblright{}.) Say it once more.
\end{tcolorbox}
\section[la désinformation]{la désinformation — disinformation, false information}

\label{lesson:FR-C285-desinformation}



\begin{quote}
Before the new one: say the French for the rumour, then the French for the scoop.
\end{quote}

\subsection*{You'll want to know: la désinformation}
\textbf{la désinformation} — \textquotedblleft{}disinformation, false information\textquotedblright{}.

\begin{quote}
Il y a beaucoup de désinformation sur internet.
\end{quote}

\begin{quote}
\emph{There is a lot of false information on the internet.}
\end{quote}

People also say \emph{les fausses nouvelles}, or \emph{les fake news} in speech.

\begin{cousinweb}
\emph{désinformation} is built on \emph{information}, a learned borrowing from Latin \emph{īnfōrmātiō} \textquotedblleft{}an outline, an idea\textquotedblright{}, built on \emph{fōrma} \textquotedblleft{}shape\textquotedblright{} (Lewis \& Short).
\end{cousinweb}

\begin{grammarlens}[title={What you've built}]
One more word for telling, explaining and describing at B1.
\end{grammarlens}

\subsection*{Your turn}
\begin{itemize}
\raggedright
  \item \emph{Say it:} \emph{la désinformation}
  \item \emph{Say it:} \emph{la désinformation}, once more
  \item \emph{Say it:} a sentence of your own with \emph{la désinformation}
  \item \emph{Recall:} say the French for the rumour, then the French for the scoop, then say \emph{la désinformation} again
\end{itemize}

\begin{tcolorbox}[breakable,colback=teal!4,colframe=teal!35!black,title={Before you move on}]
What does la désinformation mean? (\textquotedblleft{}Disinformation, false information\textquotedblright{}.) Say it once more.
\end{tcolorbox}
\section[la censure]{la censure — censorship}

\label{lesson:FR-C285-censure}



\begin{quote}
Before the new one: say the French for the scoop, then the French for disinformation.
\end{quote}

\subsection*{You'll want to know: la censure}
\textbf{la censure} — \textquotedblleft{}censorship\textquotedblright{}.

\begin{quote}
Le film a été interdit par la censure.
\end{quote}

\begin{quote}
\emph{The film was banned by the censors.}
\end{quote}

The verb is \emph{censurer}.

\begin{cousinweb}
\emph{censure} is a learned borrowing from Latin \emph{cēnsūra} \textquotedblleft{}the office of the censor\textquotedblright{}, from \emph{cēnsēre} \textquotedblleft{}to assess, to judge\textquotedblright{} (Lewis \& Short).
\end{cousinweb}

\begin{grammarlens}[title={What you've built}]
One more word for telling, explaining and describing at B1.
\end{grammarlens}

\subsection*{Your turn}
\begin{itemize}
\raggedright
  \item \emph{Say it:} \emph{la censure}
  \item \emph{Say it:} \emph{la censure}, once more
  \item \emph{Say it:} a sentence of your own with \emph{la censure}
  \item \emph{Recall:} say the French for the scoop, then the French for disinformation, then say \emph{la censure} again
\end{itemize}

\begin{tcolorbox}[breakable,colback=teal!4,colframe=teal!35!black,title={Before you move on}]
What does la censure mean? (\textquotedblleft{}Censorship\textquotedblright{}.) Say it once more.
\end{tcolorbox}
\section[la polémique]{la polémique — the controversy, the row}

\label{lesson:FR-C285-polemique}



\begin{quote}
Before the new one: say the French for disinformation, then the French for censorship.
\end{quote}

\subsection*{You'll want to know: la polémique}
\textbf{la polémique} — \textquotedblleft{}the controversy, the row\textquotedblright{}.

\begin{quote}
Ses paroles ont créé une polémique.
\end{quote}

\begin{quote}
\emph{His words caused a row.}
\end{quote}

As an adjective: \emph{un ton polémique}.

\begin{cousinweb}
\emph{polémique} comes from Greek \emph{polemikós} \textquotedblleft{}warlike\textquotedblright{} (EtymDB), built on \emph{pólemos} \textquotedblleft{}war\textquotedblright{}: a war of words.
\end{cousinweb}

\begin{grammarlens}[title={What you've built}]
One more word for telling, explaining and describing at B1.
\end{grammarlens}

\subsection*{Your turn}
\begin{itemize}
\raggedright
  \item \emph{Say it:} \emph{la polémique}
  \item \emph{Say it:} \emph{la polémique}, once more
  \item \emph{Say it:} a sentence of your own with \emph{la polémique}
  \item \emph{Recall:} say the French for disinformation, then the French for censorship, then say \emph{la polémique} again
\end{itemize}

\begin{tcolorbox}[breakable,colback=teal!4,colframe=teal!35!black,title={Before you move on}]
What does la polémique mean? (\textquotedblleft{}The controversy, the row\textquotedblright{}.) Say it once more.
\end{tcolorbox}
